\documentclass[12pt]{article}
\usepackage[utf8]{inputenc}
\usepackage[T1]{fontenc}
\usepackage{hyperref}
\title{Tahitian Language Resources for K-12 Curriculum Development}
\author{Compiled by Claude Sonnet 5 (claude-sonnet-5) for Andr\'as Kornai}
\date{}
\begin{document}
\maketitle

\section{Introduction}
Tahitian (reo Tahiti; ISO 639-3: tah; Glottocode: tahi1242) is an Austronesian language and an official language of French Polynesia. The Acad\'emie tahitienne (Fare V\=ana'a), founded in 1972, standardizes vocabulary, grammar and spelling. Resources are far richer than for the other languages in this batch, though this search was brief.

\section{Primary Educational Resources}
\begin{itemize}
\item The Acad\'emie has produced school textbooks, reported as \emph{Ta'u puta reo Tahiti} (primary) and \emph{Hei pua ri'i} (secondary), per the Wikipedia summary: \url{https://en.wikipedia.org/wiki/Tahitian_Academy}
\item Official French education document noting a weekly Tahitian hour in 6\`eme: \url{https://eduscol.education.fr/document/1043/download}
\end{itemize}

\section{Linguistic and Reference Materials}
\begin{itemize}
\item Acad\'emie tahitienne publications (download pages): \url{https://www.farevanaa.pf/fra/download/71} and \url{https://www.farevanaa.pf/fra/download/3}
\item Fare V\=ana'a website news (Government service-public.pf): \url{https://www.service-public.pf/dgen/evolution-du-site-internet-du-fare-vanaa/}
\item Lexilogos Tahitian dictionary links: \url{https://www.lexilogos.com/tahitien_dictionnaire.htm}
\end{itemize}

\section{Religious and Community Resources}
None specifically identified in search results.

\section{Digital and Audio Resources}
\begin{itemize}
\item Reo app and site, a bilingual Tahitian-French tool based on a 17,000-entry dictionary from the Acad\'emie: \url{https://outremers360.com/culture/polynesie-une-version-trilingue-pour-lapplication-et-le-site-reo}
\item Learn Tahitian terms (dance-oriented, informal): \url{https://tahitidanceonline.com/learn-tahitian-terms-with-reo/}
\end{itemize}

\section{Curriculum Development Context}
Tahitian is taught in schools in French Polynesia and the University of French Polynesia offers training in Polynesian languages. Textbooks exist for primary and secondary levels, but a full K-12 sequence taught in Tahitian was not documented in this search. Next steps: contact the Acad\'emie tahitienne and the local education directorate.
\end{document}